\scanpage{104}
\[
u:\underbrace{\widehat A^{**}}_{\Vert\ \prod A^{ij}}\unsure{-|\longrightarrow}\underbrace{\widehat {A'}{}^{**}}_{\Vert\ \prod A'{}^{ij}}
\]
hom d'Anneaux différentiels appliquant $A^{**+}$ dans $A'{}^{**+}$ et compatible aux p.d., et tels que $u(T)=T$ \unsure{NB} a priori
\[
u:\underbrace{k\{T\}}_{\Vert\ H^*(A^{**})}\longrightarrow\underbrace{k\{T\}}_{\Vert\ H^*(A'{}^{**})}
\]
hom compatible aux p.d.). \deleted{À vrai dire, il \unsure{se peut}\ \unsure{déduire de ce qui précède}}

On désigne par
\[
\mathrm{DR}_{\mathrm{sp}}'(X,k)_0,\quad\mathrm{DR}_{\mathrm{spufl}}'(X,k)_0
\]
les cat. dérivées associées (i.e. les groupoïdes engendrés). On a un \deleted{\unsure{des}} diagramme comm. de foncteurs
\diagram{\begin{tikzcd}[row sep=large,column sep=large]
\mathrm{DR}_{\mathrm{sp}}(X,k)\arrow[r,"{\text{surj. sur obj}}"]\arrow[d,"{\deleted{\unsure{iso.}}\ \text{surjectif sur obj}}"]&\mathrm{DR}_{\mathrm{sp}}(X,k)_0\simeq\text{ponctuelle}\arrow[d]\\
\mathrm{DR}_{\mathrm{sp}}'(X,k)\arrow[r,"{\text{surj. sur objet}}"']&\mathrm{DR}_{\mathrm{sp}}'(X,k)_0\\
\mathrm{DR}_{\mathrm{spufl}}(X,k)\arrow[r,"{\text{surj. sur obj}}"]\arrow[uu,bend left=25,"{\deleted{\unsure{\downarrow}}}"']\arrow[d,"{\text{surjectif sur obj}}"]&\mathrm{DR}_{\mathrm{spufl}}(X,k)_0\simeq\text{ponct.}\arrow[uu,bend left=25,"{\unsure{\uparrow|\downarrow}}"]\arrow[d,"{\text{(si ufl. stable par les }\Gamma^i,\ i>0\text{)}}"]\\
\mathrm{DR}_{\mathrm{spufl}}'(X,k)\arrow[r,"{\text{surj. sur ob}}"]\arrow[uu,bend left=25,"{\deleted{\unsure{\downarrow}}}"']&\mathrm{DR}_{\mathrm{spufl}}'(X,k)_0\simeq\arrow[uu,bend left=25,"{\deleted{\unsure{\downarrow}}}"]
\end{tikzcd}}
qui montre que les deux catégories \inserted{\unsure{dér.}} $\mathrm{DR}_{\mathrm{sp}}'(X,k)_0$ et $\mathrm{DR}_{\mathrm{spufl}}'(X,k)_0$ sont des groupoïdes \inserted{\unsure{ens.}} \unsure{ponctuels|pointés} (munis d'une foule d'objets avec syst. transitif d'iso. entre). Donc ils sont \unsure{can.}

\margin{103}
